\scanpage{217}
\mspage{111}

Considérons l'application

\[
\begin{aligned}
F:\ \Pi\times S&\longrightarrow S^\vee\\
(\varphi,s)&\longmapsto\varphi(s)
\end{aligned}
\]

on aura donc

\[
\underbrace{F(T_\lambda\varphi,s)}_{(\varphi\circ\lambda^{-1})(s)=\varphi(\lambda^{-1}s}
=F(\varphi,\lambda_S^{-1}s)
\]

ce qui montre que $F$ se factorise en une application canonique, qui est en fait iso d'ens.

\[
\Pi\wedge_Z(S!)\xrightarrow{\sim} S^\vee\qquad\varphi\wedge s\longmapsto\varphi(s) \tag*{(76)}
\]

où $S!$ désigne le $Z$-ensemble déduit de $S$ en y remplaçant l' \deleted{\unsure{structure|action}} \inserted{opération} donnée \inserted{``à gauche''} de $Z$ par l'opération ``à droite'' que s'en déduit par $\lambda\longmapsto\lambda^{-1}$.

\deleted{Comme les applications $\varphi$ commutent aux opérations de $Z$ sur $S$ et $S^\vee$} \inserted{Dans cet isomorphisme, on aura}

\[
\deleted{\unsure{et}}\quad
\underbrace{\lambda.(\varphi\wedge s)}_{\substack{(T_\lambda\varphi)\wedge s=\varphi\wedge\lambda_S^{-1}s\\\text{par définition}}}
\longrightarrow\lambda_{S^\vee}(\varphi(s))
\]

on voit que (76) est un isomorphisme de $Z$-ensembles. \deleted{\unsure{C}} Est-ce même un isomorphisme de $\unsure{\mathcal G|G}$-\deleted{torseurs} \inserted{ensembles} ? (En notant que du fait que $\mathcal G$ agit \deleted{\unsure{comme gr}} \inserted{sur $S!$} par $Z$-automorphismes, il agit aussi sur l'ensemble tordu $\Pi\wedge_Z(S!)$ par transport de structure, par

\[
\begin{aligned}
&g(\varphi\wedge s)=\varphi\wedge g_S(s)\qquad)\\
&\text{On aura}\quad g(\varphi\wedge s)=\varphi\wedge g_S(s)\longmapsto
\varphi(g_S(s))\overset?=g_{S^\vee}\varphi(s).
\end{aligned} \tag*{(77)}
\]

% archive sheet number: 216
